% -*- coding: utf-8 -*-
% Latex基础语法.tex
% Latex基础语法
\documentclass[UTF8]{ctexart}
\usepackage{minted}
\usepackage{geometry}

\geometry{a4paper,scale=0.8}

\usepackage{amsmath}
\usepackage{float}  % 提供 [H] 选项
\usepackage[svgnames]{xcolor}
\usepackage{tcolorbox}
\tcbuselibrary{listings,skins}

% 定义 latexexample 环境
\newtcblisting{latexexample}[1][]{
    colback=gray!10,
    colframe=gray!50,
    listing side text,  % 左侧代码，右侧效果
    title={LaTeX 示例},
    fonttitle=\bfseries,
    #1
}

\title{\LaTeX 基础语法}
\author{埃及猪肉}
\date{\today}



\begin{document}
\maketitle
\tableofcontents


\section{简单结构}
\begin{minted}[linenos, frame=single]{latex}
\documentclass[UTF8]{ctexart}

\usepackage{amsmath}
\usepackage{float}  % 提供 [H] 选项
\usepackage[svgnames]{xcolor}
\usepackage{geometry}
\geometry{a4paper, scale=0.8}

\title{文档标题}
\author{作者姓名}
\date{\today}

\begin{document}

\maketitle

\section{引言}

这是一个简单的段落。LaTeX会自动处理段落缩进和换行。

\subsection{子标题}

这是子标题下的段落内容。

\end{document}
\end{minted}

\section{文字排版}

\subsection{字符输入}

\subsubsection{换行和换页}

在\LaTeX{} 中不会自动换行，需要间隔一行，或者手动输入换行符\verb|\\|、\verb|\par|或换行命令\verb|\newline|。


\begin{latexexample}%
我是第一行。
我是第二行？

我是第二行。\\我是第三行。

我是第四行。\newline 我是第五行。

我是第六行。\par 我是第七行。

我是第八行。

我是第九行。
\end{latexexample}

使用\verb|\newline|和\verb|\\|进行换行时，行首会顶格；而使用\verb|\par|和间隔一行进行换行时，行首会缩进。


\subsubsection{字体}

字体有多种样式，如加粗、斜体、下划线、等宽字体、无衬线字体等。\par
相应的命令为：\verb|\textbf{}|、\verb|\textit{}|、\verb|\underline{}|、\verb|\texttt{}|、\verb|\textsf{}|。

\begin{latexexample}%
\textbf{加粗}

\textit{斜体}

\underline{下划线}

\texttt{等宽字体}

\textsf{无衬线字体}

\textsf{\textbf{加粗无衬线字体}}
\end{latexexample}


\subsubsection{字体大小}

可以通过\verb|\tiny|、\verb|\scriptsize|、\verb|\footnotesize|、\verb|\small|、\verb|\normalsize|、\verb|\large|、\verb|\Large|、\verb|\LARGE|、\verb|\huge|、\verb|\Huge|命令设置字体大小。

\begin{latexexample}%
正常字体

\tiny{小号字体}

\scriptsize{更小号字体}

\footnotesize{脚注大小}

\small{小号字体}

\normalsize{正常大小}

\large{大号字体}

\Large{更大号字体}

\LARGE{更大号字体}

\huge{超大号字体}

\Huge{超大号字体}
\end{latexexample}


\subsubsection{字体颜色和背景色}

字体颜色可以通过\verb|\textcolor{}{}|命令设置。

\begin{latexexample}%
不加颜色

\textcolor{red}{红色字体}

\textcolor{blue}{蓝色字体}

\textcolor{green}{绿色字体}

\textcolor{orange}{橙色字体}

\textcolor{purple}{紫色字体}

\textcolor{cyan}{青色字体}

\textcolor{yellow}{黄色字体}

\textcolor{gray}{灰色字体}

\textcolor{brown}{棕色字体}

\textcolor{magenta}{品红色字体}

\textcolor{olive}{橄榄色字体}

\end{latexexample}

颜色可以通过颜色代码来设置颜色，也可以使用RGB编码/十六进制颜色码。

\begin{latexexample}%
不加代码

\textcolor{red}{红色字体}

\textcolor[rgb]{1,0,0}{RGB颜色字体}

\end{latexexample}

背景色可以通过\verb|\colorbox{}{}|命令设置。

\begin{latexexample}%
不加颜色

\colorbox{red}{红色背景}

\colorbox{blue}{蓝色背景}

\colorbox{green}{绿色背景}

\colorbox{orange}{橙色背景}

\colorbox{purple}{紫色背景}

\colorbox{cyan}{青色背景}

\colorbox{yellow}{黄色背景}

\colorbox{gray}{灰色背景}

\end{latexexample}




\subsubsection{公式加框}

可以使用\verb|\boxed{}|、\verb|\fbox{}|、\verb|\framebox{}|命令加框。

其中\verb|boxed{}|是amsmath的数学模式命令，必须用\verb|$...$|包裹。

\begin{latexexample}%
$\boxed{\text{我是加框的文字1}}$

\fbox{我是加框的文字2}

\framebox{我是加框的文字3}
\end{latexexample}



\subsubsection{特殊符号的输入}

\LaTeX 中,一般字符可以被正常输入\par
而一些字符如\verb|&|、\verb|#|、\verb|%|、\verb|{ }|、\verb|_|、\verb|^|、\verb|~|、\verb|\|等在LaTeX中有特殊的意义，如果需要输入它们本身，需要使用转义字符\verb|\|。

\begin{latexexample}%
输出如下：
\&
\#
\%
\{ \}
\_
\^{}
\~{}
\textbackslash
\end{latexexample}

以上代码分别输出：\& \# \% \{ \} \_ \^{} \~{} \textbackslash

\subsubsection{水平间距}

\LaTeX{} 中，可以使用\verb|\hspace{}|命令来加入空格并设置水平间距。\par
单位有cm、mm、pt、ex、em、in等，含义如下：

\begin{itemize}
    \item cm：厘米
    \item mm：毫米
    \item pt：磅
    \item ex：字体大小
    \item em：字体大小
    \item in：英寸
\end{itemize}

\begin{latexexample}%
水平距离为1cm\hspace{1cm}的间距。

水平距离为1mm\hspace{1mm}的间距。

水平距离为1pt\hspace{1pt}的间距。

水平距离为1ex\hspace{1ex}的间距。

水平距离为1em\hspace{1em}的间距。

水平距离为1in\hspace{1in}的间距。

\end{latexexample}

同样的，还可以配合下划线命令\verb|\underline{}|来显示空格的效果：\par
\begin{latexexample}%
水平距离为1cm\underline{\hspace{1cm}}的间距。

水平距离为1mm\underline{\hspace{1mm}}的间距。

水平距离为1pt\underline{\hspace{1pt}}的间距。

水平距离为1ex\underline{\hspace{1ex}}的间距。

水平距离为1em\underline{\hspace{1em}}的间距。

水平距离为1in\underline{\hspace{1in}}的间距。

\end{latexexample}

\subsubsection{列表环境}

\LaTeX 中提供了多种列表环境，包括无序列表、有序列表和描述列表。

无序列表使用\verb|itemize|环境：

\begin{minted}[frame=single]{latex}
\begin{itemize}
\item 无序列表项
\item 无序列表项
\item 无序列表项
\end{itemize}
\end{minted}

效果如下：
\begin{itemize}
\item 无序列表项
\item 无序列表项
\item 无序列表项
\end{itemize}

有序列表使用\verb|enumerate|环境：

\begin{minted}[frame=single]{latex}
\begin{enumerate}
\item 有序列表项
\item 有序列表项
\item 有序列表项
\end{enumerate}
\end{minted}

效果如下：
\begin{enumerate}
\item 有序列表项
\item 有序列表项
\item 有序列表项
\end{enumerate}

描述列表使用\verb|description|环境：

\begin{minted}[frame=single]{latex}
\begin{description}
\item[描述列表项1] 描述列表项1的内容。
\item[描述列表项2] 描述列表项2的内容。
\item[描述列表项3] 描述列表项3的内容。
\end{description}
\end{minted}

效果如下：
\begin{description}
\item[描述列表项1] 描述列表项1的内容。
\item[描述列表项2] 描述列表项2的内容。
\item[描述列表项3] 描述列表项3的内容。
\end{description}




\subsection{数学输入}

\subsubsection{基本元素}

\paragraph{希腊字母}\mbox{}

希腊字母可以通过反斜杠加字母名称输入。

\begin{table}[H]
\centering
\caption{小写希腊字母}
\begin{tabular}{cc|cc|cc}
\hline
代码 & 输出 & 代码 & 输出 & 代码 & 输出 \\
\hline
\verb|$\alpha$|   & $\alpha$   & \verb|$\iota$|    & $\iota$    & \verb|$\sigma$|   & $\sigma$   \\
\verb|$\beta$|    & $\beta$    & \verb|$\kappa$|   & $\kappa$   & \verb|$\tau$|     & $\tau$     \\
\verb|$\gamma$|   & $\gamma$   & \verb|$\lambda$|  & $\lambda$  & \verb|$\upsilon$| & $\upsilon$ \\
\verb|$\delta$|   & $\delta$   & \verb|$\mu$|      & $\mu$      & \verb|$\phi$|     & $\phi$     \\
\verb|$\epsilon$| & $\epsilon$ & \verb|$\nu$|      & $\nu$      & \verb|$\chi$|     & $\chi$     \\
\verb|$\zeta$|    & $\zeta$    & \verb|$\xi$|      & $\xi$      & \verb|$\psi$|     & $\psi$     \\
\verb|$\eta$|     & $\eta$     & \verb|$\pi$|      & $\pi$      & \verb|$\omega$|   & $\omega$   \\
\verb|$\theta$|   & $\theta$   & \verb|$\rho$|     & $\rho$     &                    &            \\
\hline
\end{tabular}
\end{table}

\begin{table}[H]
\centering
\caption{大写希腊字母}
\begin{tabular}{cc|cc|cc}
\hline
代码 & 输出 & 代码 & 输出 & 代码 & 输出 \\
\hline
\verb|$\Gamma$|   & $\Gamma$   & \verb|$\Xi$|    & $\Xi$    & \verb|$\Phi$|   & $\Phi$   \\
\verb|$\Delta$|   & $\Delta$   & \verb|$\Pi$|    & $\Pi$    & \verb|$\Psi$|   & $\Psi$   \\
\verb|$\Theta$|   & $\Theta$   & \verb|$\Sigma$| & $\Sigma$ & \verb|$\Omega$| & $\Omega$ \\
\verb|$\Lambda$|  & $\Lambda$  & \verb|$\Upsilon$| & $\Upsilon$ &               &          \\
\hline
\end{tabular}
\end{table}

\begin{table}[H]
\centering
\caption{希腊字母变体}
\begin{tabular}{cc|cc}
\hline
代码 & 输出 & 代码 & 输出 \\
\hline
\verb|$\varepsilon$| & $\varepsilon$ & \verb|$\varrho$|    & $\varrho$    \\
\verb|$\vartheta$|  & $\vartheta$  & \verb|$\varsigma$|  & $\varsigma$  \\
\verb|$\varpi$|     & $\varpi$     & \verb|$\varphi$|    & $\varphi$    \\
\hline
\end{tabular}
\end{table}


\paragraph{上标和下标}\mbox{}

指数或者上标可以通过\verb|^|来输入；
下标可以通过\verb|_|来输入。

上标和下标默认只能作用于一个字符；
如果需要作用于多个字符，可以使用大括号\verb|{}|来包含多于一个字符的内容。

\begin{latexexample}%
$a^2$

$a_2$

$a^22$

$a^{22}$

$a_22$

$a_{22}$

\end{latexexample}

\paragraph{分数和根式}\mbox{}

分数可以通过\verb|\frac{}{}|来输入。

\begin{latexexample}%
$\frac{a+b}{c+d}$
\end{latexexample}

根式可以通过\verb|\sqrt[]{}|来输入。

\begin{latexexample}%
$\sqrt[1]{2}$
\end{latexexample}

\paragraph{运算符}\mbox{}

常用的运算符有：\verb|+|、\verb|-|、\verb|\times|、\verb|\div|、\verb|\pm|、\verb|\mp|、\verb|\cdot|等。

\begin{latexexample}%
$a+b$

$a-b$

$a\times b$

$a\div b$

$a\pm b$

$a\mp b$

$a\cdot b$
\end{latexexample}



\end{document}
